\subsection{最初に必要な用語}
\par\smallskip\noindent\refstepcounter{table}\label{tab:ch09-01}表 \thetable: 最初に必要な用語における用語・初学者向けの説明\par\smallskip
表 \thetable は、最初に必要な用語における用語・初学者向けの説明を比較・確認しやすい形で整理したものである。\par\smallskip
\begin{longtable}{>{\raggedright\arraybackslash}p{0.25\linewidth}>{\raggedright\arraybackslash}p{0.70\linewidth}}
\toprule
用語 & 初学者向けの説明 \\
\midrule
\endfirsthead
\toprule
用語 & 初学者向けの説明 \\
\midrule
\endhead
ソフトウェアエンジニアリング管理 & ソフトウェアプロジェクトを計画し、見積り、測定し、監視し、制御し、調整し、リスクを管理し、チームを導く活動である。英語では Software Engineering Management、略して SEM である。 \\
プロジェクト管理 & 期限、費用、範囲、品質、資源、リスク、関係者などを扱い、プロジェクト目標を達成するための管理である。 \\
\term{測定}{Measurement} & 成果物、プロセス、資源に値やラベルを割り当て、判断に使える情報を作ることである。単に数値を集めるだけではなく、管理判断に結び付ける必要がある。 \\
ソフトウェア開発ライフサイクル & 構想から開発、運用、保守、終了までの流れである。英語では Software Development Life Cycle、略して SDLC である。 \\
予測型 & 要求や設計を早い段階で詳しく定め、詳細な計画に基づいて進めるライフサイクルである。 \\
適応型 & 短い反復を通じて要求や計画を段階的に詳しくし、変化へ対応するライフサイクルである。 \\
\term{成果物}{Deliverable} & 仕様書、設計書、検査報告、テスト済みソフトウェアなど、測定可能で確認可能な仕事の産物である。 \\
作業分解構造 & 作業を小さな単位に分解した構造である。英語では Work Breakdown Structure、略して WBS である。 \\
リスク & 目標達成に影響する不確実な事象である。悪い影響だけでなく、機会として良い影響を持つ場合もある。 \\
\term{情報製品}{Information Product} & 測定データを分析し、判断に使える形にしたグラフ、指標、報告、結論である。 \\
\bottomrule
\end{longtable}
